\begin{abstract}
Retrieval-derived uncertainty can carry error information, but it also depends on which reference documents are selected. We conducted a development-informed follow-up of a fixed 17-class clinical section-routing Cascade, using 773 reference, 258 development and 257 previously reported historical-test records from MTS-Dialog. With prediction probabilities fixed, matched-count filtered and unfiltered evidence sets differed for 103 development records, and retrieval-label entropy changed for 93. In grouped cross-fitting over six assignments of the same development records, retrieval features improved ranking metrics, but evidence-view benefits depended on the metric and coverage. On the historical test partition, which was previously reported and is not an untouched holdout, a combined selector fitted on development data had a discrete AURC of 0.104713 versus 0.109681 for a predictive-only selector (difference $-0.004967$; conditional 95\% paired group-bootstrap interval $-0.018524$ to $0.007669$). The point estimates favored the combined selector, but the interval included zero, and the combined selector retained one error at 25\% coverage where the predictive-only selector retained none. Parallel analyses with three alternative predictors (logistic regression, XGBoost and a calibrated linear SVM on TF-IDF features) showed no distinguishable gain from the same retrieval features (AURC differences $-0.0015$, $-0.0007$ and $-0.0046$; all 95\% intervals included zero). We therefore do not claim a superior selector, a formal risk guarantee or external generalization. The study reuses data and predictions from a preceding study.
\end{abstract}
